\documentclass{deimj}

\usepackage[T1]{fontenc}
\usepackage{lmodern}

% 必要に応じてパッケージを導入して下さい。
\usepackage[dvipdfm]{graphicx}
%\usepackage{latexsym}
%\usepackage{txfonts}
%\usepackage[fleqn]{amsmath}
%\usepackage[psamsfonts]{amssymb}
%\usepackage[deluxe]{otf}

% 印刷位置調整 %
% 必要に応じて値を変更してください．
\hoffset -10mm % <-- 左に 10mm 移動
\voffset -10mm % <-- 上に 10mm 移動

\newcommand{\AmSLaTeX}{%
 $\mathcal A$\lower.4ex\hbox{$\!\mathcal M\!$}$\mathcal S$-\LaTeX}
\newcommand{\PS}{{\scshape Post\-Script}}
\def\BibTeX{{\rmfamily B\kern-.05em{\scshape i\kern-.025em b}\kern-.08em
 T\kern-.1667em\lower.7ex\hbox{E}\kern-.125em X}}

\jtitle{F1レースの状況・環境指標を用いた時系列相関量計量による \\ 戦略的タイヤ交換タイミングの提案}
% \jsubtitle{サブタイトル} % サブタイトルを付けたいときはこの行の先頭の % を取る
\authorlist{%
 \authorentry{坂野 心咲}{Misaki Sakano}{MU}
 \authorentry{林 康弘}{Yasuhiro Hayashi}{MU} 
 \authorentry{清木 康}{Yasushi Kiyoki}{MU}
}

\affiliate[MU]{武蔵野大学 データサイエンス学科\hskip1zw
  〒135--8181 東京都江東区有明 3--3--3}
 {Department of Data Science, Musashino University\\
  Ariake, Koto-ku, Tokyo 135--8181, Japan}

% 同一組織で、複数のメールアドレスのドメインがある場合、
% ダガーがずれてしまう。
% その場合、以下のコメントアウトを解除すると、
% 自由にダガーを設定可能であることが分かる。
\MailAddress{$\dagger$ s2122084@stu.musashino-u.ac.jp,
\{yhayashi, y-kiyoki\}@musashino-u.ac.jp
}


\begin{document}
\pagestyle{empty}

\begin{jabstract}
Formula1(F1)のレース戦略の重要な要素の１つは，タイヤ戦略にある．ストラテジストの形式知・暗黙知を反映した知識ベースの構築とその効果的な活用によって，特定の専門家の経験値のみに依存することなく， F1レースタイヤ戦略の最適解を導出することは重要な課題となっている．本研究は、F1を対象として，レースの状況・環境に関する指標を対象とした時系列相関量計量によるタイヤ戦略構築支援方式を提案する．本方式の特徴は，レースに勝つための最適な戦略を，過去の成功事例から導くため，レースデータをデータベース化し，時系列相関量計量によって現在の事象に対して過去の類似する事象を発見する点にある．これにより， 状況・環境が類似する過去の成功事例（周回単位）をもとに最適と考えらえるタイヤ戦略支援の実現を目指す．本研究の実現システムとして，実際のF1データを用いて専門家の知識ベース・過去成功事例データベースを構築し，Pythonを用いた実験システムを実装した．この実験システムにより，過去成功事例を反映した最適なタイヤ戦略の提案を実現でき，本方式の実現可能性を検証した
\end{jabstract}

\begin{jkeyword}
時系列相関量計量、F1レースデータベース、戦略的タイヤ交換タイミング、過去との類似性、成功事例検索
\end{jkeyword}

\maketitle

\section{はじめに}

あ

\subsection{研究背景}

あ

\subsection{現状の課題}

あ

\subsection{本研究の目的}

あ

\section{提案方式}

あ

\subsection{全体概要}

あ

\subsection{F1レースデータベースの構築}

あ

\subsection{データの前処理と特徴量抽出}

あ

\subsection{時系列相関量計量による類似事例検索}

あ

\section{実装}

本章では，第2章で提案した方式に基づき，実際に構築した類似事例検索システムの実装について述べる．
本実装では，分析の目的に応じて，以下の2つの検索方法を実装した．
\begin{enumerate}
  \item 単周検索: 特定の1周をクエリとして入力し，その周回におけるレース状況・環境指標が類似する過去の事例を検索する方法
  \item 複数周検索: タイヤ交換の前後を含む連続した周回を対象として，その間のラップタイムの推移やレース状況・環境指標の変化を含めた俯瞰的な比較を行い，類似する過去の事例を検索する方法
\end{enumerate}

\subsection{実装環境}

本システムの実装は，Google Colaboratory上で行った．使用言語はPython3.12であり，数値計算やデータ分析には以下のライブラリを利用した．

\begin{itemize}
  \item \textbf{pandas}: 時系列データの操作，データフレーム形式での集計処理
  \item \textbf{NumPy}: 数値計算処理，全体的なデータ操作
  \item \textbf{SciPy}: ユークリッド距離の算出(\texttt{cdist})，および回帰分析による傾きの算出(\texttt{linregress})
  \item \textbf{scikit-learn}: 特徴量の正規化(Min-Maxスケーリング等)の前処理
\end{itemize}
開発環境としてGoogle Colaboratoryを使用した理由は，クラウド上で実行環境が統一されているため再現性が高く，大規模なF1データの処理に必要な計算リソースの確保やライブラリ変更が柔軟に行えるためである．

\subsection{データ取得と前処理の実装}

第2章で述べたF1レースデータベースの構築において，実際のF1レースデータを国際自動車連盟（FIA）の公式サイトや公式中継映像等から収集し，各レースと各ドライバー毎にCSVファイルを作成した．作成したCSVファイルをGoogle Colaboratory上にアップロードし，Pandasライブラリを用いてデータフレームとして読み込んだ．
次に，レースごとに分かれているファイルを結合するため，年を抽出し，year列を追加した上で結合処理を行った．さらに，ラップタイムとロスタイムを計算しやすくする為にミリ秒に変換し，欠損値の処理や不要な列の削除などの前処理を行い，分析に適した形式に整形した．さらに，特徴量抽出のために，レース状況・環境指標から必要な特徴量を選択する．

複数周検索においては，指定した周回間のラップタイムを用いた新たな特徴量を生成した．具体的には，指定した周回間のラップタイムの平均値，標準偏差，傾き（回帰分析による）を計算し，新たな特徴量として追加した．これにより，複数周にわたるレース状況の変化を捉えやすくし，類似事例検索の精度向上を図った．

また，前処理の重要なステップとして，正規化処理を行った．本データセットには，ラップタイム（ミリ秒単位の大きな値）や順位（1桁の小さな値）といった異なるスケールの変数が混在している．これらをそのまま用いると，値の大きな変数が距離計算の結果に大きく影響してしまうため，単周検索および複数周検索の両方において，各特徴量に対してMin-Maxスケーリングを行い，0から1の範囲に正規化した．
これにより，すべての特徴量が同等の重みで類似度算出に寄与するように調整した．

\subsection{類似事例検索ロジックの実装}

類似度の算出には，ユーグリッド距離を使用し，実装においては，SciPyライブラリの\texttt{cdist}関数を用いた．
単周検索では，クエリとして指定された1周分の特徴量ベクトルとデータベース内の全周回の特徴量ベクトルとの距離を計算し，距離が最小となる上位5件の事例を抽出した．
複数周検索では，対象区間と過去データのスライディングウィンドウ間の特徴量ベクトルを用いて同様に距離計算を行い，その総和（または平均）を区間全体の類似度として定義した．これにより，指定された複数周にわたるレース状況の変化を考慮した類似事例検索を実現した．

\subsection{出力結果の形式}

検索結果は，計算された類似度の高い順に上位5件の事例を表示する形で実装した．出力形式は，分析者が直感的に状況を把握できるよう表形式とした．各事例については，以下の情報を含めた表形式で出力した．

すべての検索モードにおいて共通して出力される情報は以下の通りである．
\begin{itemize}
  \item \textbf{基本情報}: 開催年，周回数，ドライバーID
  \item \textbf{類似度}: 算出されたユークリッド距離（0に近いほど類似している）
  \item \textbf{レース状況}: 順位(Rank)，タイヤの使用周回数(UseLap)，タイヤ種別(Tire)，天候指標(Weather)
\end{itemize}

さらに，複数周検索においては，指定された区間の傾向を可視化するため，以下の統計量も併せて出力する．
\begin{itemize}
  \item \textbf{LapTime Mean}: 区間内の平均ラップタイム
  \item \textbf{LapTime Std}: ラップタイムの標準偏差（タイムのばらつき）
  \item \textbf{LapTime Slope}: ラップタイムの回帰直線の傾き（デグラデーションの進行度合い）
\end{itemize}

また，複数周検索では，検索の対象にした年度を除いた過去データからの検索結果も合わせて表示するようにした．これにより，現在のレース状況に類似する過去の成功事例を効果的に参照できるようにした．

\begin{table}[htbp]
  \centering
  \caption{出力結果の例（単周検索モード）}
  \label{tab:result_single}
  \resizebox{\columnwidth}{!}{% ページ幅に合わせて縮小
  \begin{tabular}{ccccccccc}
    \hline
    Year & Lap & Distance & Rank & LapTime & LossTime & UseLap & Tire & Weather \\
    \hline
    2022 & 26 & 0.087 & 8 & 152713 & 0 & 6 & 2 & 1 \\
    2022 & 24 & 0.111 & 8 & 134670 & 0 & 4 & 2 & 1 \\
    2022 & 23 & 0.307 & 8 & 96720 & 0 & 3 & 2 & 1 \\
    2022 & 27 & 0.349 & 8 & 86589 & 0 & 7 & 2 & 1 \\
    2022 & 28 & 0.401 & 8 & 85828 & 0 & 8 & 2 & 1 \\
    \hline
  \end{tabular}
  }
\end{table}


\begin{table}[htbp]
  \centering
  \caption{出力結果の例（複数周検索モード・別年度抽出）}
  \label{tab:result_multi}
  \resizebox{\columnwidth}{!}{% ページ幅に合わせて縮小
  \begin{tabular}{ccccccccc}
    \hline
    Year & Lap & Driver & Dist & Time Mean & Time Std & Slope & Tire & Weath \\
    \hline
    2023 & 14 & 3 & 0.419 & 98002 & 33562 & -4003 & 2 & 1 \\
    2023 & 17 & 3 & 0.421 & 86617 & 7895 & 380 & 2 & 1 \\
    2023 & 15 & 3 & 0.487 & 87652 & 7744 & 680 & 2 & 1 \\
    2023 & 18 & 3 & 0.493 & 86515 & 7943 & 45 & 2 & 1 \\
    2023 & 16 & 3 & 0.527 & 86916 & 7786 & 639 & 2 & 1 \\
    \hline
  \end{tabular}
  }
\end{table}


\section{過去データを用いた評価シミュレーション（計画）}

本提案手法の有効性を検証するため，2023年の実際のレースデータを用いたバックテスト（シミュレーション）を行いたいと考えている．本セクションでは，その評価計画について述べる．

\subsection{実験の目的}

本実験の目的は，提案方式が過去の類似事例に基づいて算出する「戦略的推奨」が，実際のレースにおける「正解（勝者の行動）」とどの程度一致するかを検証することである．
「過去成功事例を反映した戦略」の有効性を確認するため，抽出された類似過去事例群が当時どのような戦略（タイヤ選択・ピットタイミング）を採用して成功したかを集計し，それを実際のレースに対する推奨行動として確率的に提示できるかをシミュレーションする．

\subsection{実験設定と使用データ}

評価には，2021年〜2024年のF1公式データ（FastF1より取得）を使用する．
具体的には，戦略判断が重要となった過去のレース（例：天候変化時や接戦時）を「テストシナリオ」として選定する．
各シナリオにおいて，戦略決定が行われた直前の周回までのデータをシステムに入力し，類似する過去の事例を抽出する．

\subsection{評価手順}

本実験では，過去のレース結果を用いたバックテスト（振り返り検証）形式で評価を行う．具体的な手順は以下の通りである．

\begin{enumerate}
  \item \textbf{テストケースの選定}:
  評価対象とするレース（例：2023年 オーストラリアGP）と，戦略判断を行う時点（例：20周目）を決定する．この時点より後のデータ（21周目以降の結果）は，システムには入力せず「未知の情報」として扱う．

  \item \textbf{類似事例の検索}:
  決定した時点までのレースデータをクエリとしてシステムに入力し，類似度が最も高い過去の事例を$N$件（本実験では5件とする）抽出する．

  \item \textbf{推奨戦略の集計（確率的投票）}:
  抽出された5件の類似事例において，そのドライバーがその後どのような行動をとって成功したかを集計する．
  （例：5件中4件が「翌周にタイヤ交換」を行っていた場合，本システムの推奨戦略は「確率80\%でタイヤ交換」となる）

  \item \textbf{実データとの答え合わせ}:
  システムが導き出した推奨戦略と，実際のテスト対象レースにおける勝者（または上位入賞者）の行動を比較する．
  システムの推奨（例：タイヤ交換）と，実際の勝者の行動（例：タイヤ交換して優勝）が一致した場合を「正解（有効な提案）」と判定する．
\end{enumerate}

\subsection{評価指標}

提案方式の実現可能性を以下の指標で定量的に評価する．
\begin{itemize}
  \item \textbf{戦略一致率}: システムが推奨した戦略（例：80\%の類似事例がピットインを選択）と，実際のレース勝者の行動が一致した割合．
  \item \textbf{事例の妥当性評価}: 抽出された類似事例が，レース解説やレポート等の定性的評価と照らし合わせて，本当に「状況が似ていた」と言えるかどうかのケーススタディ．
\end{itemize}

\section{おわりに}

\subsection{まとめ}

本研究では，F1レースの状況・環境指標を用いた時系列相関量計量による戦略的タイヤ交換タイミングの提案方式を提案した．本方式は，過去の成功事例データベースを構築し，現在のレース状況に類似する過去の事例を時系列相関量計量によって検索することで，最適なタイヤ戦略を支援するものである．実装システムを通じて，本方式の実現可能性を検証する．今後は，提案方式の有効性を評価するために，過去データを用いた評価シミュレーションを実施する予定である．

\subsection{今後の展望}

今後の課題として，提案方式の精度向上と実用化に向けた検討が挙げられる．具体的には，より多様なレース状況に対応できるよう，データベースの拡充と特徴量抽出手法の改善を図る必要がある．また，仮定に基づいた計算を行っているが，この仮定が正しいのかの検証を行うことも重要である．さらに，実際のレース環境での適用可能性を検証し，ストラテジストへのフィードバックループを構築することも重要である．これにより，F1レースにおける戦略的タイヤ交換タイミングの最適化に寄与できると期待される．

\section*{謝辞}
謝辞を記載する場合はこちら．

%\vspace{30mm} <- 文献が本文と近すぎるときは適宜利用してください．
\vspace{2em}

\begin{thebibliography}{99}
\bibitem{Codd1970}
  E. F. Codd, 
  ``A Relational Model of Data for Large Shared Data Banks,''
  Communications of the {ACM} (CACM), Vol. 13, No. 6, pp. 377--387, 1970.
\end{thebibliography}

% bibtexを利用する場合の例
% \bibliographystyle{jplain}
% \bibliography{reference}

\end{document}
